\documentclass[leqno]{article}
% Shared mark scheme preamble for resources.danieljsmith.org: the pack
% preamble, unchanged, then the scheme's own layer. A scheme's .tex is
% rendered from its YAML (scripts/README.md, Mark schemes), never edited.
\input{\CurrentFilePath/pack_preamble.tex}

% The page: equal side margins, so the table runs the full width of the
% header rule. The footers' right edge follows the text's.
\geometry{left=0.8in, right=0.8in}
\makeatletter
\renewcommand{\djs@footR}[1]{%
  \put(\dimexpr1in+\hoffset+\oddsidemargin+\textwidth\relax,
    -\dimexpr\paperheight-0.32in\relax){\makebox[0pt][r]{#1}}}
\makeatother

\usepackage[export]{adjustbox}
\usepackage{tabularray}
\definecolor{msGrid}{RGB}{176, 172, 164}
\definecolor{msTint}{RGB}{245, 243, 239}

% The header on every page, the front page included, and the booklet label
% under the contents-page title. \djsQuestionPageStyle stays empty.
\newcommand{\djsSchemeHeader}[1]{%
  \djsPackHeader[markscheme]{\textit{Mark scheme} - #1}%
  \renewcommand{\djsPackTitle}{#1}%
  \renewcommand{\djsBookletLabel}{Mark scheme}%
}

% The front page: the contents heading, spaced as tocloft spaces the packs'
% "Questions", then the contents table. A question's name and page number
% both link to its first page.
\newcommand{\msContentsHeading}{%
  \vskip\cftbeforetoctitleskip
  {\noindent\cfttoctitlefont\contentsname\par}%
  \vskip\cftaftertoctitleskip}
\NewTblrEnviron{mscontents}
\SetTblrInner[mscontents]{
  colspec = {Q[wd=26mm, l] X[l] Q[wd=16mm, c] Q[wd=14mm, r]},
  width = \linewidth,
  rowsep = 4pt, colsep = 6pt,
  hlines = {0.4pt, msGrid}, vlines = {0.4pt, msGrid},
  hline{2} = {2.2pt, scarlet},
  row{1} = {font=\small\bfseries},
  row{Z} = {font=\bfseries},
  column{2} = {font=\small},
}
\newcommand{\msContentsQuestion}[1]{\hyperref[q:#1]{Question #1}}
\newcommand{\msContentsPage}[1]{\hyperref[q:#1]{\pageref*{q:#1}}}
\newcommand{\msSep}{\hspace{0.9em}}

% The conventions pages' layout: headings in the documents' own bold, tables
% ruled and padded as the scheme's, roomy lists.
\newcommand{\msHead}[1]{%
  \par\vspace{18pt}%
  {\noindent\bfseries #1\par}%
  \vspace{6pt}}
\NewTblrEnviron{mstblr}
\SetTblrInner[mstblr]{rowsep = 6pt, colsep = 6pt,
  hlines = {0.4pt, msGrid}, vlines = {0.4pt, msGrid}}
\newenvironment{msList}
  {\begin{itemize}[leftmargin=1.2em, itemsep=5pt, topsep=3pt, parsep=0pt]}
  {\end{itemize}}

% A question's first page: the link target and the PDF bookmark, as
% \questionitem gives the packs.
\newcommand{\msQnum}{}
\newcommand{\msQuestion}[1]{%
  \clearpage
  \renewcommand{\msQnum}{#1}%
  \phantomsection\label{q:#1}%
  \addcontentsline{toc}{subsection}{\protect\djsTocEntry{#1}}}

% The scheme's table: Question 16mm, Scheme the rest, Marks 16mm, Notes
% 44mm; every cell ruled in msGrid; the scarlet rule under the head row. A
% long table has no caption or "continued" foot; a continuation page opens
% with a muted line above the repeated head row. The rules across a part's
% rows, the alternative's tint and the total row are the rendered table's.
\NewTblrTheme{ms}{%
  \DefTblrTemplate{firsthead}{default}{}%
  \DefTblrTemplate{middlehead,lasthead}{default}{%
    {\small\itshape\color{djsMuted}Question \msQnum\ continued\par}%
    \vspace{6pt}}%
  \DefTblrTemplate{firstfoot,middlefoot,lastfoot}{default}{}%
}
\NewTblrEnviron{msscheme}
\SetTblrOuter[msscheme]{long, theme=ms, entry=none, label=none}
\SetTblrInner[msscheme]{
  colspec = {Q[wd=16mm, l] X[l] Q[wd=16mm, c] Q[wd=44mm, l]},
  width = \linewidth,
  rowsep = 6pt, colsep = 6pt,
  rowhead = 1,
  vlines = {0.4pt, msGrid},
  hline{1} = {0.4pt, msGrid}, hline{2} = {2.2pt, scarlet},
  hline{Z} = {0.4pt, msGrid},
  column{1} = {font=\bfseries}, column{3} = {font=\bfseries},
  column{4} = {font=\small},
  row{1} = {rowsep=4pt, font=\small\bfseries},
}

% Two marks on a row: each note starts with its code.
\newcommand{\msNote}[2]{\textbf{#1}\enspace #2}
% A full-row figure: the worked solution's own picture, scaled uniformly and
% only down, to the table less the Question column. The picture is drawn
% with the line width it has in the solutions document, at its list depth
% there (the optional argument, 1 by default: a part's working), so one
% sized by \linewidth keeps its size; tabularray sets the cell's \linewidth
% to the cell's measured width, which would shrink it at every pass. The
% widths are measured once, in a solution box at the packs' text width
% (pack_preamble.tex's margins, 0.8in and 1in).
\newlength{\msSpanW}
\makeatletter
\newcommand{\ms@keepLW}[1]{%
  \expandafter\xdef\csname ms@solnLW@#1\endcsname{\the\linewidth}}
\AtBeginDocument{%
  \setlength{\msSpanW}{\dimexpr\textwidth-16mm-28pt\relax}%
  \setbox\z@\vbox{%
    \begin{minipage}{\dimexpr\paperwidth-1.8in\relax}
      \begin{djsSolution}
        \ms@keepLW{0}%
        \begin{enumerate}\item \ms@keepLW{1}%
        \begin{enumerate}\item \ms@keepLW{2}%
        \begin{enumerate}\item \ms@keepLW{3}%
        \end{enumerate}\end{enumerate}\end{enumerate}
      \end{djsSolution}
    \end{minipage}}}
\newcommand{\msFigure}[2][1]{%
  \adjustbox{max width=\msSpanW, max totalheight=0.6\textheight, valign=t}{%
    \hsize=\csname ms@solnLW@#1\endcsname
    \linewidth=\hsize
    #2}}
\makeatother
% A figure mark's feature, in the Scheme column.
\newcommand{\msFeature}[1]{\textit{#1}}
% An alternative's title row.
\newcommand{\msAlternative}[1]{\textit{Alternative:} #1}
% Drafts only: a queried mark's code, the query's why under its note, and a
% flag in its part's place.
\newcommand{\msQueried}[2]{#1~\(\to\) #2?}
\newcommand{\msQuery}[1]{{\itshape\color{djsMuted}Query: #1}}
\newcommand{\msFlag}[2]{\textbf{#1:}\enspace #2}


\djsSchemeHeader{Maclaurin Series}
\djsSitePath{Core Pure}{Further Calculus}

\djsLastUpdated{25 October 2026}
\begin{document}
\djsSetQuestionCount{17}
\djsAddMarks{1}{6}
\djsAddMarks{2}{6}
\djsAddMarks{3}{8}
\djsAddMarks{4}{6}
\djsAddMarks{5}{7}
\djsAddMarks{6}{6}
\djsAddMarks{7}{6}
\djsAddMarks{8}{10}
\djsAddMarks{9}{5}
\djsAddMarks{10}{8}
\djsAddMarks{11}{5}
\djsAddMarks{12}{6}
\djsAddMarks{13}{12}
\djsAddMarks{14}{10}
\djsAddMarks{15}{8}
\djsAddMarks{16}{9}
\djsAddMarks{17}{10}
\vspace*{6pt}
\djsContentsTitle
\vspace{30pt}
\msContentsHeading
\begin{mscontents}{}
Question & Marks by part & Marks & Page \\
\msContentsQuestion{1} & (a)(i)\,2\msSep (a)(ii)\,2\msSep (b)\,2 & 6 & \msContentsPage{1} \\
\msContentsQuestion{2} & (a)\,2\msSep (b)\,1\msSep (c)\,3 & 6 & \msContentsPage{2} \\
\msContentsQuestion{3} & (i)\,5\msSep (ii)\,3 & 8 & \msContentsPage{3} \\
\msContentsQuestion{4} & (a)(i)\,2\msSep (a)(ii)\,2\msSep (a)(iii)\,2 & 6 & \msContentsPage{4} \\
\msContentsQuestion{5} & (a)\,2\msSep (b)\,2\msSep (c)\,3 & 7 & \msContentsPage{5} \\
\msContentsQuestion{6} & (a)\,2\msSep (b)\,1\msSep (c)\,3 & 6 & \msContentsPage{6} \\
\msContentsQuestion{7} & (a)\,2\msSep (b)\,2\msSep (c)\,2 & 6 & \msContentsPage{7} \\
\msContentsQuestion{8} & (a)\,5\msSep (b)\,3\msSep (c)\,1\msSep (d)\,1 & 10 & \msContentsPage{8} \\
\msContentsQuestion{9} & & 5 & \msContentsPage{9} \\
\msContentsQuestion{10} & (i)\,4\msSep (ii)\,4 & 8 & \msContentsPage{10} \\
\msContentsQuestion{11} & & 5 & \msContentsPage{11} \\
\msContentsQuestion{12} & (a)\,3\msSep (b)\,3 & 6 & \msContentsPage{12} \\
\msContentsQuestion{13} & (i)\,3\msSep (ii)\,4\msSep (iii)\,5 & 12 & \msContentsPage{13} \\
\msContentsQuestion{14} & (a)(i)\,4\msSep (a)(ii)\,2\msSep (b)\,4 & 10 & \msContentsPage{14} \\
\msContentsQuestion{15} & (a)\,4\msSep (b)\,4 & 8 & \msContentsPage{15} \\
\msContentsQuestion{16} & (a)\,4\msSep (b)\,5 & 9 & \msContentsPage{16} \\
\msContentsQuestion{17} & (a)(i)\,4\msSep (a)(ii)\,2\msSep (b)\,4 & 10 & \msContentsPage{17} \\
\SetCell[c=2]{r} Total & & 128 & \\
\end{mscontents}
\newpage
% The conventions pages, shared by every pack's mark scheme: one source of
% truth, \input after the front page, with no strand or pack logic. The keys
% on the first side, the rules on the second. The scheme layer defines
% \msHead, msList and mstblr.
\vspace*{6pt}%
{\noindent\cfttoctitlefont Reading the mark scheme\par}%
\vspace{10pt}%
These rules apply throughout, unless a note says otherwise.

\msHead{Marks}
Every mark is worth one, so \textbf{M1 A1} on a row is two marks. Only marks
coded ft follow through.\par
\vspace{10pt}%
\begin{mstblr}{
  colspec = {Q[wd=15mm, l] X[l]},
  width = \linewidth,
  hline{2} = {2.2pt, scarlet},
  column{1} = {font=\bfseries},
  row{1} = {font=\small\bfseries, rowsep=4pt},
}
Code & Meaning \\
M1 & Method: a correct method, applied to this question. A slip doesn't cost it if the method is clear. A formula must be used, not just quoted. \\
dM1 & Dependent method: as M1, but only if the method mark before it was earned, unless its note names another. \\
A1 & Accuracy: a correct answer or step. It needs the method mark it follows. \\
A1ft & As A1, following through from an earlier answer that the note names. \\
A1~AG & The last mark of a part whose answer is printed in the question: complete working, with no errors. \\
B1 & A correct value, fact or statement. It needs no method mark. \\
B1ft & As B1, following through from an earlier answer that the note names. \\
E1 & A written explanation, interpretation or conclusion. Concluding a proof by induction needs every earlier mark. \\
\end{mstblr}

\msHead{Abbreviations}
ft and AG qualify a code in the Marks column. The others appear in the
notes.\par
\vspace{10pt}%
\begin{mstblr}{
  colspec = {Q[wd=17mm, l] X[l] Q[wd=15mm, l] X[l]},
  width = \linewidth,
  column{1} = {font=\bfseries}, column{3} = {font=\bfseries},
}
ft & follow through & AG & answer given in the question \\
cao & correct answer only, with no follow-through & cso & correct solution only: no errors in the part \\
oe & or an equivalent form & awrt & any answer that rounds to this value \\
soi & seen, or implied by later working & isw & ignore working after a correct answer \\
SC & special case: credit outside the main method & their & the student's own earlier value \\
condone & a slip that doesn't cost the mark & [\;\,] & working the mark doesn't need \\
s.f. & significant figures & d.p. & decimal places \\
\end{mstblr}

\newpage
\msHead{Reading the scheme}
\begin{msList}
\item Each row is a line of working, with the marks it earns beside it.
  Where a row carries two marks, each note starts with its code.
\item \textbf{[\textit{n}]} ends each part with its total.
\item A shaded \textit{Alternative} is a second method for the same total.
  Mark an answer by the method it follows.
\item Any other valid method earns the equivalent marks, unless the question
  names a method or says \textit{hence}.
\item A sketch earns one B1 for each feature named below its figure.
\end{msList}

\msHead{Accuracy and answers}
\begin{msList}
\item Unless a question says otherwise, give answers to 3 significant
  figures, and angles in degrees to 1 decimal place. Angles in radians take
  3 significant figures.
\item A final answer with more figures is accepted if it rounds to the
  scheme's value; one with fewer is not. An answer spoiled by rounding too
  early loses its accuracy mark.
\item Where a question asks for an exact answer, a decimal loses the final
  accuracy mark. Otherwise either is accepted.
\item Wrong or missing units cost nothing unless the question asks for them.
\item A correct answer with no working earns full marks, unless the question
  needs the working: show that, prove, hence, a method it names, or an exact
  answer. A wrong answer with no working earns nothing. Calculator answers
  are treated the same way.
\item A correct answer, once reached, isn't spoiled by later working (isw).
\end{msList}

\msHead{Value of $\boldsymbol{g}$}
\begin{msList}
\item Take $g = 9.8\,\mathrm{m\,s^{-2}}$ unless a question states a value,
  which is then the only one accepted. Answers that use $g = 9.8$ may be
  given to 2 or 3 significant figures.
\item Where a question doesn't state $g$, an answer worked with $g = 10$
  throughout earns the same marks, keeping any value printed in the
  question. The worked solutions list the final answers with $g = 10$.
\end{msList}


\msQuestion{1}
\begin{msscheme}{
  hline{3-4,6-7,9-10} = {2-Z}{0.4pt, msGrid},
  hline{5,8,11} = {0.4pt, msGrid},
  row{11} = {font=\bfseries},
}
Question & Scheme & Marks & Notes \\
1(a)(i) & {\(f'(x) = -2e^{-2x}\)} & M1 & {Differentiates to the form \(ke^{-2x}\), any nonzero constant \(k\).} \\
& {\(f''(x) = 4e^{-2x}\)} & A1 & {Both derivatives correct.} \\
\nopagebreak & & [2] & \\
1(a)(ii) & {\(f(0) = 1,\ f'(0) = \text{"}{-2}\text{"},\ f''(0) = \text{"}4\text{"}\)\\[4pt] \(e^{-2x} = 1 + x(\text{"}{-2}\text{"}) + \frac{x^2}{2!}(\text{"}4\text{"}) + \cdots\)} & M1 & {Uses \(f(0) + xf'(0) + \frac{x^2}{2!}f''(0)\) with their derivatives from (a)(i) at \(x = 0\). Not for replacing \(x\) by \(-2x\) in the series for \(e^x\).} \\
& {\(e^{-2x} = 1 - 2x + 2x^2 + \cdots\)} & A1 & \\
\nopagebreak & & [2] & \\
1(b) & {\(x = \tfrac18\)\\[4pt] \(e^{-1/4} \approx 1 - 2\left(\tfrac18\right) + 2\left(\tfrac18\right)^2\)} & M1 & {Substitutes \(x = \frac18\) into their three-term series from (a)(ii).} \\
& {\(\phantom{e^{-1/4}} = 1 - \tfrac14 + \tfrac1{32} = \tfrac{25}{32}\)} & A1~AG & {From the correct series, reaching \(\frac{25}{32}\) without error; combined arithmetic accepted. Not for a calculator value of \(e^{-1/4}\).} \\
\nopagebreak & & [2] & \\
\nopagebreak \SetCell[c=2]{r} Total for Question 1 & & 6 & \\
\end{msscheme}

\msQuestion{2}
\begin{msscheme}{
  hline{3-4,6,8-10} = {2-Z}{0.4pt, msGrid},
  hline{5,7,11} = {0.4pt, msGrid},
  row{11} = {font=\bfseries},
}
Question & Scheme & Marks & Notes \\
2(a) & {\(e^x = 1 + x + \frac{x^2}{2!} + \frac{x^3}{3!} + \frac{x^4}{4!} + \cdots\)} & M1 & {Uses the standard series, or derivatives at \(x = 0\) in the Maclaurin formula, with factorial denominators. Need not be simplified.} \\
& {\(\phantom{e^x} = 1 + x + \frac{x^2}{2} + \frac{x^3}{6} + \frac{x^4}{24}\)} & A1 & {All five terms, with the factorials evaluated.} \\
\nopagebreak & & [2] & \\
2(b) & {\(e^{-x} = 1 - x + \frac{x^2}{2} - \frac{x^3}{6} + \frac{x^4}{24}\)} & B1ft & {ft their five terms in (a) with \(x\) replaced by \(-x\), signs simplified.} \\
\nopagebreak & & [1] & \\
2(c) & {\(\cosh x = \tfrac12\left(e^x + e^{-x}\right)\)\\[4pt] \(\phantom{\cosh x} = 1 + \frac{x^2}{2} + \frac{x^4}{24} + \cdots\)} & M1 & {Adds their series for \(e^x\) and \(e^{-x}\) and halves, reaching at least the terms in \(x^2\) and \(x^4\).} \\
& {\(\cosh^2\left(\tfrac{x}{2}\right) = \tfrac12\left(\text{"}1 + \frac{x^2}{2} + \frac{x^4}{24}\text{"} + 1\right)\)} & dM1 & {Adds 1 to their series for \(\cosh x\) and divides by 2.} \\
& {\(\phantom{\cosh^2\left(\tfrac{x}{2}\right)} = 1 + \tfrac14 x^2 + \tfrac1{48}x^4 + \cdots\)} & A1 & {Correct series, or \(a = 1\), \(b = \frac14\), \(c = \frac1{48}\) stated.} \\
\nopagebreak & & [3] & \\
\nopagebreak \SetCell[c=2]{r} Total for Question 2 & & 6 & \\
\end{msscheme}

\msQuestion{3}
\begin{msscheme}{
  hline{3-6,8-10} = {2-Z}{0.4pt, msGrid},
  hline{7,11} = {0.4pt, msGrid},
  row{11} = {font=\bfseries},
}
Question & Scheme & Marks & Notes \\
3(i) & {\(\sin 2x = 2x - \frac{(2x)^3}{3!} + \frac{(2x)^5}{5!} + \cdots\)\\[4pt] \(\cos 3x = 1 - \frac{(3x)^2}{2!} + \frac{(3x)^4}{4!} + \cdots\)} & M1 & {Replaces \(x\) by \(2x\) in the \(\sin x\) series and by \(3x\) in the \(\cos x\) series, at least two terms each. Need not be simplified.} \\
& {\(\cos x - \cos 3x = \left(1 - \frac{x^2}{2} + \frac{x^4}{24}\right)\)\\[4pt] \(\phantom{\cos x - \cos 3x =} - \text{"}\left(1 - \frac{9x^2}{2} + \frac{81x^4}{24}\right)\text{"}\)\\[4pt] \(\phantom{\cos x - \cos 3x} = 4x^2 - \frac{10}{3}x^4 + \cdots\)} & M1 & {Subtracts their \(\cos 3x\) series from the \(\cos x\) series, reaching terms in \(x^2\) and \(x^4\). Condone sign slips.} \\
& {\(\text{"}\left(2x - \frac43x^3 + \frac{4}{15}x^5\right)\text{"}\,\text{"}\left(4x^2 - \frac{10}{3}x^4\right)\text{"}\)\\[4pt] \(\phantom{A} = 8x^3 - \frac{20}{3}x^5 - \frac{16}{3}x^5 + \cdots\)} & dM1 & {Multiplies their two series, forming the \(x^3\) term and both products giving \(x^5\). Condone slips in coefficients.} \\
& {\(\sin 2x(\cos x - \cos 3x) = 8x^3 - 12x^5 + \cdots\)} & A1 A1 & {\msNote{A1}{\(8x^3\)}\\ \msNote{A1}{\(-12x^5\)}} \\
\nopagebreak & & [5] & \\
3(ii) & {\(2\left(\text{"}8x^3 - 12x^5\text{"}\right) = x\)} & M1 & {Replaces the LHS by twice their series from (i), of the form \(ax^3 + bx^5\), and equates to \(x\).} \\
& {\(24x^4 - 16x^2 + 1 = 0\)\\[4pt] \(\Rightarrow x^2 = \frac{16 \pm \sqrt{256 - 96}}{48} = \ldots\)} & dM1 & {Divides by \(x\) and solves their three-term quadratic in \(x^2\) by the formula with their values or completing the square.} \\
& {\(x^2 = \frac{4 - \sqrt{10}}{12} \Rightarrow x = \sqrt{\frac{4 - \sqrt{10}}{12}}\)} & A1~AG & {From a correct quadratic in \(x^2\) without error, selecting the smaller value of \(x^2\) and the positive root.} \\
\nopagebreak & & [3] & \\
\nopagebreak \SetCell[c=2]{r} Total for Question 3 & & 8 & \\
\end{msscheme}

\msQuestion{4}
\begin{msscheme}{
  hline{3-4,6-7,9-10} = {2-Z}{0.4pt, msGrid},
  hline{5,8,11} = {0.4pt, msGrid},
  row{11} = {font=\bfseries},
}
Question & Scheme & Marks & Notes \\
4(a)(i) & {\(f'(x) = \frac{\cos x}{1 + \sin x}\)\\[4pt] \(f''(x) = \frac{(-\sin x)(1 + \sin x) - \cos x \cos x}{(1 + \sin x)^2}\)} & M1 & {Obtains \(f'(x)\) of the form \(\frac{k\cos x}{1 + \sin x}\) and differentiates it by the quotient or product rule. Condone sign slips.} \\
& {\(f''(x) = -\frac{1}{1 + \sin x}\)} & A1 & {Any correct form, simplified or not.} \\
\nopagebreak & & [2] & \\
4(a)(ii) & {\(f(0) = 0,\ f'(0) = 1,\ f''(0) = \text{"}{-1}\text{"}\)\\[4pt] \(f(x) = 0 + x(1) + \frac{x^2}{2!}(\text{"}{-1}\text{"}) + \cdots\)} & M1 & {Substitutes \(x = 0\) into their \(f'\) and \(f''\) and uses \(f(0) + xf'(0) + \frac{x^2}{2!}f''(0)\), or composes the series for \(\ln(1+u)\) and \(\sin x\).} \\
& {\(f(x) = x - \frac{x^2}{2} + \cdots\)} & A1 & \\
\nopagebreak & & [2] & \\
4(a)(iii) & {\(\int_0^{2/5} \text{"}x - \frac{x^2}{2}\text{"}\,\mathrm{d}x = \left[\frac{x^2}{2} - \frac{x^3}{6}\right]_0^{2/5}\)\\[4pt] \(\phantom{\int_0^{2/5}} = \frac{(2/5)^2}{2} - \frac{(2/5)^3}{6}\)} & M1 & {Integrates their two-term series, raising the power of at least one term, and substitutes \(x = \frac25\).} \\
& {\(= \frac{26}{375} \approx 0.069333\)} & A1 & {awrt 0.069333; \(\frac{26}{375}\) alone or fewer than 6 d.p. scores A0.} \\
\nopagebreak & & [2] & \\
\nopagebreak \SetCell[c=2]{r} Total for Question 4 & & 6 & \\
\end{msscheme}

\msQuestion{5}
\begin{msscheme}{
  hline{3-4,6-7,9-11} = {2-Z}{0.4pt, msGrid},
  hline{5,8,12} = {0.4pt, msGrid},
  row{12} = {font=\bfseries},
}
Question & Scheme & Marks & Notes \\
5(a) & {\(\ln\left(\frac{1+2x}{1-x}\right) = \ln(1+2x) - \ln(1-x)\)\\[4pt] \(\ln(1+2x) = 2x - \frac{(2x)^2}{2} + \frac{(2x)^3}{3} + \cdots\)\\[4pt] \(\ln(1-x) = -x - \frac{x^2}{2} - \frac{x^3}{3} - \cdots\)} & M1 & {Splits the logarithm and substitutes \(2x\) and \(-x\) into the \(\ln(1+x)\) series, to at least the \(x^3\) terms. Condone sign slips.} \\
& {\(\left(2x - 2x^2 + \frac83x^3\right) - \left(-x - \frac12x^2 - \frac13x^3\right)\)\\[4pt] \(\phantom{A} = 3x - \frac32x^2 + 3x^3\)} & A1~AG & {Correct expansions of both logarithms combined without error to the printed terms.} \\
\nopagebreak & & [2] & \\
5(b) & {Series valid only for \(-\frac12 < x \leq \frac12\)} & B1 & {Accept \(|2x| < 1\) or \(|x| < \frac12\); condone the endpoints.} \\
& {\(x = \frac58 > \frac12\), outside the range of validity} & E1 & {States that \(\frac58\) lies outside a correct range, so the series is not valid there. Not for saying only that the answer is inaccurate.} \\
\nopagebreak & & [2] & \\
5(c) & {\(\ln 2 \approx 3\left(\tfrac14\right) - \tfrac32\left(\tfrac14\right)^2 + 3\left(\tfrac14\right)^3\)\\[4pt] \(\ln 8 = 3\ln 2 \approx 3 \times \ldots\)} & M1 & {Substitutes \(x = \frac14\) into the three terms as an approximation for \(\ln 2\) and multiplies by 3. Condone arithmetic slips.} \\
& {\(\ln 8 \approx \frac{135}{64} \approx 2.11\)} & A1 & {\(\frac{135}{64}\) oe, or awrt 2.11.} \\
& {\(x = \tfrac14 \Rightarrow \frac{1+2x}{1-x} = \frac{3/2}{3/4} = 2\)\\[4pt] \(-\tfrac12 < \tfrac14 \leq \tfrac12\), within the range of validity} & E1 & {Shows \(\frac{1+2x}{1-x} = 2\) at \(x = \frac14\) and states that \(\frac14\) lies in the range of validity.} \\
\nopagebreak & & [3] & \\
\nopagebreak \SetCell[c=2]{r} Total for Question 5 & & 7 & \\
\end{msscheme}

\msQuestion{6}
\begin{msscheme}{
  hline{3-4,6,8-10} = {2-Z}{0.4pt, msGrid},
  hline{5,7,11} = {0.4pt, msGrid},
  row{11} = {font=\bfseries},
}
Question & Scheme & Marks & Notes \\
6(a) & {\(-\ln(1-x) = x + \frac{x^2}{2} + \frac{x^3}{3} + \cdots\)} & M1 & {Writes the series for \(\ln(1-x)\) or \(-\ln(1-x)\) with at least two terms and an indication it continues. Condone a sign slip.} \\
& {\(0 < x < 1 \Rightarrow \frac{x^2}{2} + \frac{x^3}{3} + \cdots > 0\)\\[4pt] \(\phantom{0 < x < 1} \Rightarrow -\ln(1-x) > x\)} & A1~AG & {Correct series, with the terms after \(x\) stated positive for \(0 < x < 1\), leading to the printed inequality. A0 without the reason.} \\
\nopagebreak & & [2] & \\
6(b) & {\(x = 1 - \tfrac1t\), \(t > 1 \Rightarrow 0 < x < 1\)\\[4pt] \(-\ln\left(\tfrac1t\right) > 1 - \tfrac1t \Rightarrow \ln t > 1 - \tfrac1t\)} & B1 & {Substitutes \(x = 1 - \frac1t\) into (a), states \(0 < x < 1\) for \(t > 1\), and reaches the printed result without error.} \\
\nopagebreak & & [1] & \\
6(c) & {\(\ln\left(\tfrac54\right) > 1 - \frac{1}{5/4}\)} & M1 & {Substitutes \(t = \frac54\) into the inequality from (b). Comparing calculator values scores nothing in this part.} \\
& {\(\ln\left(\tfrac54\right) > \tfrac15\)} & A1 & \\
& {\(5\ln\left(\tfrac54\right) > 1 = \ln e\)\\[4pt] \(\phantom{5} \Rightarrow \left(\tfrac54\right)^5 > e\)} & A1 & {cso. Multiplies by 5 and compares with \(\ln e\), or exponentiates, concluding \(\left(\frac54\right)^5\) is greater.} \\
\nopagebreak & & [3] & \\
\nopagebreak \SetCell[c=2]{r} Total for Question 6 & & 6 & \\
\end{msscheme}

\msQuestion{7}
\begin{msscheme}{
  hline{3-4,6-7,9-10} = {2-Z}{0.4pt, msGrid},
  hline{5,8,11} = {0.4pt, msGrid},
  row{11} = {font=\bfseries},
}
Question & Scheme & Marks & Notes \\
7(a) & {\(\ln(1+2x) - \ln(1-x)\)\\[4pt] \(\phantom{A} = \left(2x - \frac{(2x)^2}{2} + \frac{(2x)^3}{3}\right)\)\\[4pt] \(\phantom{A = } - \left(-x - \frac{x^2}{2} - \frac{x^3}{3}\right)\)} & M1 & {Splits the logarithm as a difference and expands both with \(\ln(1+u)\), \(u = 2x\) and \(u = -x\), at least to \(x^2\). Condone sign slips.} \\
& {\(3x - \frac32 x^2 + 3x^3\)} & A1 & {oe; isw.} \\
\nopagebreak & & [2] & \\
7(b) & {\(\frac{1+2x}{1-x} = \frac{10}{7} \Rightarrow x = \frac18\)\\[4pt] \(\ln\left(\frac{10}{7}\right) \approx 3\left(\frac18\right) - \frac32\left(\frac18\right)^2\)\\[4pt] \(\phantom{\ln\left(\frac{10}{7}\right) \approx} + 3\left(\frac18\right)^3\)} & M1 & {Solves \(\frac{1+2x}{1-x} = \frac{10}{7}\) for \(x\) and substitutes their \(x\) into their three terms from (a).} \\
& {\(\phantom{A} = \frac{192}{512} - \frac{12}{512} + \frac{3}{512} = \frac{183}{512}\)} & A1~AG & {Correct substitution of \(x = \frac18\) into correct terms, reaching \(\frac{183}{512}\) without error. Not for decimal comparison alone.} \\
\nopagebreak & & [2] & \\
7(c) & {\(\frac{1+2x}{1-x} = 10 \Rightarrow x = \frac34\)} & B1 & \\
& {\(\ln(1+2x)\) series valid only for \(-\frac12 < x \leq \frac12\)\\[4pt] \(x = \frac34\) is outside this range} & E1 & {States \(x = \frac34\) lies outside the range of validity, with \(|x| < \frac12\), \(x \leq \frac12\) or \(-\frac12 < x \leq \frac12\) seen. Not for "\(x\) too large" alone.} \\
\nopagebreak & & [2] & \\
\nopagebreak \SetCell[c=2]{r} Total for Question 7 & & 6 & \\
\end{msscheme}

\msQuestion{8}
\begin{msscheme}{
  hline{3-6,8-10,12,14} = {2-Z}{0.4pt, msGrid},
  hline{7,11,13,15} = {0.4pt, msGrid},
  row{15} = {font=\bfseries},
}
Question & Scheme & Marks & Notes \\
8(a) & {\(\sin\left(\frac{x}{2}\right) = \frac{x}{2} - \frac{x^3}{48} + \frac{x^5}{3840} + \cdots\)\\[4pt] \(\cos\left(\frac{x}{3}\right) = 1 - \frac{x^2}{18} + \frac{x^4}{1944} + \cdots\)} & M1 & {Replaces \(x\) by \(\frac{x}{2}\) in the \(\sin\) series and by \(\frac{x}{3}\) in the \(\cos\) series, at least two terms each. Need not be simplified.} \\
& {\(\frac{\sin(x/2)}{x/2} = 1 - \frac{x^2}{24} + \frac{x^4}{1920} + \cdots\)} & M1 & {Divides each term of their series for \(\sin\left(\frac{x}{2}\right)\) by \(\frac{x}{2}\). Condone one arithmetic slip.} \\
& {\(\left(1 - \frac{x^2}{24} + \frac{x^4}{1920}\right)\left(1 - \frac{x^2}{18} + \frac{x^4}{1944}\right)\)\\[4pt] \(\phantom{(} = 1 - \left(\frac{1}{24} + \frac{1}{18}\right)x^2\)\\[4pt] \(\phantom{(=} + \left(\frac{1}{1920} + \frac{1}{1944} + \frac{1}{432}\right)x^4 + \cdots\)} & dM1 & {Multiplies their two series, gathering the \(x^2\) and \(x^4\) terms. Condone one missing \(x^4\) product and arithmetic slips.} \\
& {\(= 1 - \frac{7}{72}x^2 + \frac{521}{155520}x^4\)} & A1 A1 & {\msNote{A1}{\(1 - \frac{7}{72}x^2\), simplified.}\\ \msNote{A1}{\(\frac{521}{155520}x^4\) cao; a decimal coefficient scores A0.}} \\
\nopagebreak & & [5] & \\
8(b) & {\(\int \left(\frac{1}{x} - \frac{7x}{72} + \frac{521x^3}{155520}\right)\mathrm{d}x\)\\[4pt] \(\phantom{\int} = \ln x - \frac{7x^2}{144} + \frac{521x^4}{622080}\)} & M1 & {Divides their series from (a), of the form \(1 + px^2 + qx^4\), by \(x\) and integrates, with \(\frac{1}{x} \to \ln x\) or another power increased by 1.} \\
& {\(\left[\ln x - \frac{7x^2}{144} + \frac{521x^4}{622080}\right]_{\pi/6}^{2\pi/3} = \ldots\)\\[4pt] \([= \ln 4 - \frac{35\pi^2}{1728} + \frac{8857\pi^4}{53747712}]\)} & dM1 & {Substitutes \(\frac{2\pi}{3}\) and \(\frac{\pi}{6}\) into their integral and subtracts. Implied by a correct answer.} \\
& {\(\approx 1.20244\)} & A1 & {awrt 1.20244.} \\
\nopagebreak & & [3] & \\
8(c) & {\(1.20169\)} & B1 & \\
\nopagebreak & & [1] & \\
8(d) & {\(1.20244 - 1.20169 = 0.00075\)\\[4pt] Overestimate; accurate to 3 d.p. but not to 4 d.p.} & E1 & {Makes any correct comment on the accuracy of (b), such as accurate to 3 d.p. (or 2 d.p.), or an error of 0.00075 oe. Direction not required. "Close" alone is not enough.} \\
\nopagebreak & & [1] & \\
\nopagebreak \SetCell[c=2]{r} Total for Question 8 & & 10 & \\
\end{msscheme}

\msQuestion{9}
\begin{msscheme}{
  hline{3-6} = {2-Z}{0.4pt, msGrid},
  hline{7} = {0.4pt, msGrid},
  row{7} = {font=\bfseries},
}
Question & Scheme & Marks & Notes \\
9 & {\(\frac{x}{\sin x} = \frac{1}{1 - \frac{x^2}{6} + \frac{x^4}{120} - \frac{x^6}{5040} + \cdots}\)} & M1 & {Uses the series for \(\sin x\), with at least three terms, divided by \(x\). Condone factorial or sign slips.} \\
& {\((1 + ax^2 + bx^4 + cx^6)\)\\[4pt] \(\quad\times\left(1 - \frac{x^2}{6} + \frac{x^4}{120} - \frac{x^6}{5040}\right) = 1\)} & M1 & {Multiplies by their series and sets the product equal to 1, collecting at least the terms in \(x^2\) and \(x^4\).} \\
& {\(a - \frac16 = 0\), \(b - \frac{a}{6} + \frac1{120} = 0\),\\[4pt] \(c - \frac{b}{6} + \frac{a}{120} - \frac1{5040} = 0\)\\[4pt] \(\Rightarrow a = \ldots,\ b = \ldots,\ c = \ldots\)} & dM1 & {Equates their coefficients of \(x^2\), \(x^4\) and \(x^6\) to zero and solves for all three. Condone arithmetic slips.} \\
& {\(a = \frac16\), \(b = \frac7{360}\)} & A1 & {Both exact.} \\
& {\(c = \frac{31}{15120}\)} & A1 & \\
\nopagebreak \SetCell[c=2]{r} Total for Question 9 & & 5 & \\
\end{msscheme}

\msQuestion{10}
\begin{msscheme}{
  hline{3-6,8-11} = {2-Z}{0.4pt, msGrid},
  hline{7,12} = {0.4pt, msGrid},
  row{12} = {font=\bfseries},
}
Question & Scheme & Marks & Notes \\
10(i) & {\(\ln\left(\frac{(1+x)^2}{1-2x}\right) = 2\ln(1+x) - \ln(1-2x)\)} & M1 & {Splits into multiples of \(\ln(1+x)\) and \(\ln(1-2x)\) by log laws. Condone one sign or coefficient slip.} \\
& {\(2\ln(1+x) = 2x - x^2 + \tfrac23 x^3 - \cdots\)\\[4pt] \(\ln(1-2x) = -2x - \frac{(2x)^2}{2} - \frac{(2x)^3}{3} - \cdots\)} & M1 & {Uses \(2x\) for \(x\) in the series for \(\ln(1-x)\), or \(-2x\) in that for \(\ln(1+x)\), to the \(x^3\) term. Condone sign slips.} \\
& {\(4x + x^2 + \tfrac{10}{3}x^3\)} & A1 & {All three terms correct, oe.} \\
& {\(-\tfrac12 \leq x < \tfrac12\)} & B1 & {Both endpoints, with the correct strict and non-strict inequalities.} \\
\nopagebreak & & [4] & \\
10(ii) & {\((1+x)^2 = 2(1-2x)\)\\[4pt] \(\Rightarrow x^2 + 6x - 1 = 0 \Rightarrow x = \ldots\)} & M1 & {Clears the fraction to a three-term quadratic and solves it by any method, calculator allowed. Condone sign slips.} \\
& {\(x = \sqrt{10} - 3\)} & A1 & {Or awrt 0.162, with \(-3 - \sqrt{10}\) rejected or not given.} \\
& {\(\ln 2 \approx 4(\text{"}\sqrt{10}-3\text{"}) + (\text{"}\sqrt{10}-3\text{"})^2\)\\[4pt] \(\phantom{\ln 2 \approx} + \tfrac{10}{3}(\text{"}\sqrt{10}-3\text{"})^3\)} & M1 & {Substitutes their root of the quadratic into their three-term series from (i) as the approximation to \(\ln 2\).} \\
& {\(\ln 2 \approx 0.690\)} & A1 & {0.690, to 3 d.p.} \\
\nopagebreak & & [4] & \\
\nopagebreak \SetCell[c=2]{r} Total for Question 10 & & 8 & \\
\end{msscheme}

\msQuestion{11}
\begin{msscheme}{
  hline{3-5} = {2-Z}{0.4pt, msGrid},
  hline{6} = {0.4pt, msGrid},
  row{6} = {font=\bfseries},
}
Question & Scheme & Marks & Notes \\
11 & {\(e^{x^2} = 1 + x^2 + \cdots \Rightarrow e^{x^2} - x = 1 - x + x^2 + \cdots\)} & M1 & {Replaces \(x\) by \(x^2\) in the series for \(e^x\), reaching at least \(1 + x^2\), and subtracts \(x\).} \\
& {\(e^{e^{x^2}-x} = e \times e^{u}\), \(u = -x + x^2 + \cdots\)} & M1 & {Takes out the factor \(e\), writing the expression as \(e \times e^{u}\) with \(u\) their exponent less 1.} \\
& {\(e^u = 1 + \text{"}(-x + x^2)\text{"} + \tfrac12\text{"}(-x + x^2)\text{"}^2 + \cdots\)\\[4pt] \(\phantom{e^u} = 1 - x + \ldots x^2 + \cdots\)} & dM1 & {Substitutes their \(u\), of the form \(-x + kx^2\), into \(1 + u + \frac{u^2}{2}\) and collects terms to \(x^2\). Condone slips in squaring.} \\
& {\(e^{e^{x^2}-x} = e - ex + \frac{3e}{2}x^2 + \cdots\)} & A1 A1 & {\msNote{A1}{\(e - ex\), oe such as \(e(1 - x + \cdots)\).}\\ \msNote{A1}{\(\frac{3e}{2}x^2\) oe.}} \\
\nopagebreak \SetCell[c=2]{r} Total for Question 11 & & 5 & \\
\end{msscheme}

\msQuestion{12}
\begin{msscheme}{
  hline{3-5,7-9} = {2-Z}{0.4pt, msGrid},
  hline{6,10} = {0.4pt, msGrid},
  row{10} = {font=\bfseries},
}
Question & Scheme & Marks & Notes \\
12(a) & {\(\cosh x = \tfrac12\left(e^x + e^{-x}\right)\)\\[4pt] \(\phantom{\cosh x} = 1 + \frac{x^2}{2} + \frac{x^4}{24} + \frac{x^6}{720} + \cdots\)\\[4pt] \(\sinh x = \tfrac12\left(e^x - e^{-x}\right)\)\\[4pt] \(\phantom{\sinh x} = x + \frac{x^3}{6} + \frac{x^5}{120} + \frac{x^7}{5040} + \cdots\)} & M1 & {Adds and subtracts their series for \(e^x\) and \(e^{-x}\) and halves, giving at least three terms of each series. Condone factorial slips.} \\
& {\(x\cosh x - \sinh x = \left(\tfrac12 - \tfrac16\right)x^3\)\\[4pt] \(\phantom{x\cosh x} + \left(\tfrac1{24} - \tfrac1{120}\right)x^5 + \left(\tfrac1{720} - \tfrac1{5040}\right)x^7 + \cdots\)} & dM1 & {Multiplies their \(\cosh x\) series by \(x\), subtracts their \(\sinh x\) series and collects like powers of \(x\).} \\
& {\(x\cosh x - \sinh x = \frac{x^3}{3} + \frac{x^5}{30} + \frac{x^7}{840} + \cdots\)} & A1 & \\
\nopagebreak & & [3] & \\
12(b) & {\(\frac{(ix)^3}{3} + \frac{(ix)^5}{30} + \frac{(ix)^7}{840}\)\\[4pt] \(\phantom{A} = -\frac{ix^3}{3} + \frac{ix^5}{30} - \frac{ix^7}{840}\)} & M1 & {Replaces \(x\) by \(ix\) in their three-term series from (a) and simplifies the powers of \(i\). Condone one sign slip.} \\
& {\(ix\cosh(ix) - \sinh(ix) = ix\cos x - i\sin x\)} & B1 & {Both identities used, giving \(i(x\cos x - \sin x)\) oe.} \\
& {\(x\cos x - \sin x = -\frac{x^3}{3} + \frac{x^5}{30} - \frac{x^7}{840} + \cdots\)} & A1 & {cso.} \\
\nopagebreak & & [3] & \\
\nopagebreak \SetCell[c=2]{r} Total for Question 12 & & 6 & \\
\end{msscheme}

\msQuestion{13}
\begin{msscheme}{
  hline{3-5,7-10,12-15} = {2-Z}{0.4pt, msGrid},
  hline{6,11,16} = {0.4pt, msGrid},
  row{16} = {font=\bfseries},
}
Question & Scheme & Marks & Notes \\
13(i) & {\(x = \tan y \Rightarrow 1 = \sec^2 y\,\frac{\mathrm{d}y}{\mathrm{d}x}\)} & M1 & {Writes \(x = \tan y\) and differentiates with respect to \(x\) or \(y\), reaching \(\sec^2 y\). Condone a sign slip.} \\
& {\(\frac{\mathrm{d}y}{\mathrm{d}x} = \frac{1}{\sec^2 y} = \frac{1}{1 + \tan^2 y}\)} & M1 & {Uses \(\sec^2 y = 1 + \tan^2 y\), or an equivalent identity, with \(\tan y = x\) to write the derivative in terms of \(x\).} \\
& {\(\frac{\mathrm{d}y}{\mathrm{d}x} = \frac{1}{1 + x^2}\)} & A1~AG & {Correct differentiated equation leading without error to the printed result.} \\
\nopagebreak & & [3] & \\
13(ii) & {\(\frac{1}{1 + x^2} = 1 + (-x^2) + (-x^2)^2 + \cdots\)} & M1 & {Expands \((1 + x^2)^{-1}\) by the binomial or geometric series to at least three terms. Condone sign slips.} \\
& {\(\phantom{\frac{1}{1 + x^2}} = 1 - x^2 + x^4 - \cdots\)} & A1 & \\
& {\(\arctan x = \int \text{"}1 - x^2 + x^4\text{"}\,\mathrm{d}x\)\\[4pt] \(\arctan 0 = 0 \Rightarrow c = \ldots\)} & dM1 & {Integrates their series term by term, raising each power by 1, and finds the constant from \(\arctan 0 = 0\).} \\
& {\(\arctan x = x - \frac{x^3}{3} + \frac{x^5}{5} + \cdots\)} & A1 & \\
\nopagebreak & & [4] & \\
13(iii) & {\(\ln(1 + t) = t - \frac{t^2}{2} + \frac{t^3}{3} - \frac{t^4}{4} + \cdots\)\\[4pt] with \(t = \text{"}x - \frac{x^3}{3}\text{"}\)} & M1 & {Uses the \(\ln(1 + t)\) series, with at least three terms, and \(t\) replaced by their series for \(\arctan x\).} \\
& {\(t^2 = x^2 - \frac{2}{3}x^4 + \cdots\), \(t^3 = x^3 + \cdots\),\\[4pt] \(t^4 = x^4 + \cdots\)} & M1 & {Expands at least two of \(t^2\), \(t^3\), \(t^4\) from their series, keeping terms up to \(x^4\). Condone coefficient slips.} \\
& {\(x - \frac{x^3}{3} - \frac12\left(x^2 - \frac23 x^4\right) + \frac{x^3}{3} - \frac{x^4}{4}\)} & dM1 & {Substitutes their powers into the series and collects the terms in each power up to \(x^4\).} \\
& {\(\ln(1 + \arctan x) = x - \frac{x^2}{2} + \frac{x^4}{12} + \cdots\)} & A1 A1 & {\msNote{A1}{Any two of \(x\), \(-\frac{x^2}{2}\), \(\frac{x^4}{12}\) correct.}\\ \msNote{A1}{All correct, with no \(x^3\) term other than \(0x^3\). cso.}} \\
\nopagebreak & & [5] & \\
\nopagebreak \SetCell[c=2]{r} Total for Question 13 & & 12 & \\
\end{msscheme}

\msQuestion{14}
\begin{msscheme}{
  hline{3-5,7-8,10-13} = {2-Z}{0.4pt, msGrid},
  hline{6,9,14} = {0.4pt, msGrid},
  row{14} = {font=\bfseries},
}
Question & Scheme & Marks & Notes \\
14(a)(i) & {\(f'(x) = 2\operatorname{sech}x\,(-\operatorname{sech}x\tanh x)\)\\[4pt] \(\phantom{f'(x)} = -2\operatorname{sech}^2x\tanh x\)} & M1 A1 & {\msNote{M1}{Uses the chain rule on \((\operatorname{sech}x)^2\), or a quotient rule on \(\frac{1}{\cosh^2x}\), to obtain the form \(k\operatorname{sech}^2x\tanh x\).}\\ \msNote{A1}{oe, such as \(-\frac{2\sinh x}{\cosh^3x}\).}} \\
& {\(f'' = -2f'\tanh x - 2f\operatorname{sech}^2x = 4f - 6f^2\)\\[4pt] \(f''' = 4f' - 12ff'\)\\[4pt] \(f^{(4)} = 4f'' - 12\bigl((f')^2 + ff''\bigr)\)\\[4pt] \(\Rightarrow f^{(4)}(0) = \ldots\)} & dM1 & {Differentiates three more times, in \(x\) or in terms of \(f\), and substitutes \(x = 0\). Condone sign and coefficient slips.} \\
& {\(f(0) = 1,\ f'(0) = 0,\ f''(0) = -2\)\\[4pt] \(f^{(4)}(0) = 4(-2) - 12\bigl(0 + 1(-2)\bigr)\)\\[4pt] \(\phantom{f^{(4)}(0)} = 16\)} & A1~AG & {Correct derivatives with correct values at 0 substituted, leading without error to 16.} \\
\nopagebreak & & [4] & \\
14(a)(ii) & {\(f(x) = 1 + \frac{x^2}{2!}(-2) + \frac{x^4}{4!}(16) + \cdots\)} & M1 & {Substitutes their \(f(0)\), \(f''(0)\) from (a)(i) and 16 into the Maclaurin formula with factorial denominators.} \\
& {\(\operatorname{sech}^2x = 1 - x^2 + \frac23x^4 + \cdots\)} & A1 & \\
\nopagebreak & & [2] & \\
14(b) & {\(\cos(\sqrt2\,x) = 1 - \frac{(\sqrt2\,x)^2}{2!} + \frac{(\sqrt2\,x)^4}{4!} + \cdots\)} & M1 & {Substitutes \(\sqrt2\,x\) for \(x\) in the series for \(\cos x\), reaching the \(x^4\) term.} \\
& {\(\cos(\sqrt2\,x) = 1 - x^2 + \frac16x^4 + \cdots\)} & A1 & \\
& {\(\frac{\operatorname{sech}^2x - \cos(\sqrt2\,x)}{x^4}\)\\[4pt] \(\phantom{A} = \frac{\bigl(\frac23 - \text{"}\frac16\text{"}\bigr)x^4 + \cdots}{x^4}\)} & M1 & {Subtracts their series for \(\cos(\sqrt2\,x)\) from \(1 - x^2 + \frac23x^4\) and divides by \(x^4\). Condone slips in the coefficients.} \\
& {\(\phantom{A} = \tfrac12 + \text{terms in positive powers of } x\)\\[4pt] \(\to \tfrac12\) as \(x \to 0\)} & A1~AG & {Correct series leading without error to \(\frac12\), with a statement that the terms in positive powers of \(x\) tend to 0.} \\
\nopagebreak & & [4] & \\
\nopagebreak \SetCell[c=2]{r} Total for Question 14 & & 10 & \\
\end{msscheme}

\msQuestion{15}
\begin{msscheme}{
  hline{3-6,8-11} = {2-Z}{0.4pt, msGrid},
  hline{7,12} = {0.4pt, msGrid},
  row{12} = {font=\bfseries},
}
Question & Scheme & Marks & Notes \\
15(a) & {\(2(1-\cos x) = 2\sum_{r=1}^{\infty}(-1)^{r+1}\frac{x^{2r}}{(2r)!}\)\\[4pt] \(x\sin x = \sum_{r=0}^{\infty}(-1)^r\frac{x^{2r+2}}{(2r+1)!}\)} & M1 & {Uses the general terms of the booklet's \(\cos x\) and \(\sin x\) series to write general terms for both expressions. Condone sign and index slips.} \\
& {\(x\sin x = \sum_{r=1}^{\infty}(-1)^{r-1}\frac{x^{2r}}{(2r-1)!}\)} & A1 & {Both general terms correct in powers of \(x^{2r}\), oe. Summation limits need not be stated.} \\
& {\(\frac{2(-1)^{r+1}}{(2r)!}-\frac{(-1)^{r-1}}{(2r-1)!}\)\\[4pt] \(\phantom{A} = (-1)^r\,\frac{-2+2r}{(2r)!}\)} & M1 & {Subtracts their coefficients of \(x^{2r}\) and writes them over \((2r)!\), using \((2r)! = 2r(2r-1)!\). Condone sign slips.} \\
& {\(\Rightarrow (-1)^r\frac{2(r-1)}{(2r)!}x^{2r}\)} & A1~AG & {Correct difference of coefficients leading without error to the printed general term, including \(x^{2r}\); combined steps accepted.} \\
\nopagebreak & & [4] & \\
15(b) & {\(2(1-\cos x)-x\sin x = \frac{2(2-1)}{4!}x^4+\cdots\)} & M1 & {Finds the first nonzero term of the numerator, by \(r = 2\) in (a) or by expanding the series.} \\
& {\(\phantom{2(1-\cos x)-x\sin x} = \frac{x^4}{12}+\cdots\)} & A1 & \\
& {\((1-\cos x)^2 = \left(\frac{x^2}{2}-\frac{x^4}{24}+\cdots\right)^2\)\\[4pt] \(\phantom{(1-\cos x)^2} = \frac{x^4}{4}+\cdots\)\\[4pt] \(\frac{2(1-\cos x)-x\sin x}{(1-\cos x)^2} = \frac{\text{"}\frac{x^4}{12}\text{"}+\cdots}{\frac{x^4}{4}+\cdots}\)} & M1 & {Squares their series for \(1-\cos x\) to a leading term in \(x^4\) and divides the leading terms. Condone coefficient slips.} \\
& {\(\lim_{x\to 0}\left[\ldots\right] = \frac{1/12}{1/4} = \frac13\)} & A1~AG & {Correct quotient of both series leading without error to \(\frac13\), with the limit taken.} \\
\nopagebreak & & [4] & \\
\nopagebreak \SetCell[c=2]{r} Total for Question 15 & & 8 & \\
\end{msscheme}

\msQuestion{16}
\begin{msscheme}{
  hline{3-5,7-10} = {2-Z}{0.4pt, msGrid},
  hline{6,11} = {0.4pt, msGrid},
  row{11} = {font=\bfseries},
}
Question & Scheme & Marks & Notes \\
16(a) & {\(\cosh y = \frac{x+3}{x+2}\)\\[4pt] \(\Rightarrow \sinh y\,\frac{\mathrm{d}y}{\mathrm{d}x} = \frac{(x+2)-(x+3)}{(x+2)^2}\)\\[4pt] \(\phantom{\Rightarrow \sinh y\,\frac{\mathrm{d}y}{\mathrm{d}x}} = -\frac{1}{(x+2)^2}\)} & M1 A1 & {\msNote{M1}{Differentiates \(\cosh y\) implicitly and the right side by the quotient or product rule, or differentiates arcosh by the chain rule. Condone sign slips.}\\ \msNote{A1}{Correct equation, oe; need not be simplified.}} \\
& {\(\sinh y = \sqrt{\left(\frac{x+3}{x+2}\right)^2 - 1} = \frac{\sqrt{2x+5}}{x+2}\)} & M1 & {Uses \(\sinh^2 y = \cosh^2 y - 1\) with their \(\cosh y\) to write \(\sinh y\) in terms of \(x\). Condone algebraic slips.} \\
& {\(\frac{\sqrt{2x+5}}{x+2}\,\frac{\mathrm{d}y}{\mathrm{d}x} = -\frac{1}{(x+2)^2}\)\\[4pt] \(\Rightarrow (x+2)\sqrt{5+2x}\,\frac{\mathrm{d}y}{\mathrm{d}x} + 1 = 0\)} & A1~AG & {Correct substituted equation leading without error to the printed result. Condone the positive root taken without comment.} \\
\nopagebreak & & [4] & \\
16(b) & {\((x+2)\frac{\mathrm{d}y}{\mathrm{d}x} + (5+2x)^{-\frac12} = 0\)\\[4pt] \(\Rightarrow (x+2)\frac{\mathrm{d}^2y}{\mathrm{d}x^2} + \frac{\mathrm{d}y}{\mathrm{d}x} - (5+2x)^{-\frac32} = 0\)} & M1 A1 & {\msNote{M1}{Differentiates the relation from (a), in any arrangement, or their \(\frac{\mathrm{d}y}{\mathrm{d}x}\) from it, using the product rule to obtain \(\frac{\mathrm{d}^2y}{\mathrm{d}x^2}\). Condone sign and coefficient slips.}\\ \msNote{A1}{Any correct relation or expression for \(\frac{\mathrm{d}^2y}{\mathrm{d}x^2}\); need not be simplified.}} \\
& {\(y(0) = \operatorname{arcosh}\tfrac32 = \ln\left(\tfrac32 + \sqrt{\tfrac94 - 1}\right)\)\\[4pt] \(\phantom{y(0)} = \ln\left(\frac{3+\sqrt5}{2}\right)\)} & B1 & {Needs the logarithmic form; \(\operatorname{arcosh}\frac32\) alone is not enough.} \\
& {\(y'(0) = -\frac{1}{2\sqrt5}\)\\[4pt] \(2y''(0) - \frac{1}{2\sqrt5} - \frac{1}{5\sqrt5} = 0 \Rightarrow y''(0) = \ldots\)\\[4pt] \(y = y(0) + xy'(0) + \frac{x^2}{2!}y''(0)\)} & M1 & {Substitutes \(x = 0\) into their expressions for \(y'(0)\) and \(y''(0)\) and uses them in the Maclaurin formula, including the \(\frac{1}{2!}\).} \\
& {\(a = 1,\ b = -\frac{1}{2\sqrt5},\ c = \frac{7}{40\sqrt5}\)} & A1 & {All three, oe such as \(-\frac{\sqrt5}{10}\) and \(\frac{7\sqrt5}{200}\), or awrt \(-0.224\) and \(0.0783\); or the series in the given form.} \\
\nopagebreak & & [5] & \\
\nopagebreak \SetCell[c=2]{r} Total for Question 16 & & 9 & \\
\end{msscheme}

\msQuestion{17}
\begin{msscheme}{
  hline{3-5,7-8,10-13} = {2-Z}{0.4pt, msGrid},
  hline{6,9,14} = {0.4pt, msGrid},
  row{14} = {font=\bfseries},
}
Question & Scheme & Marks & Notes \\
17(a)(i) & {\(\frac{\mathrm{d}y}{\mathrm{d}x} = n(1+\cosh x)^{n-1}\sinh x\)\\[4pt] \(\frac{\mathrm{d}^2y}{\mathrm{d}x^2} = n(n-1)(1+\cosh x)^{n-2}\sinh^2 x\)\\[4pt] \(\phantom{\frac{\mathrm{d}^2y}{\mathrm{d}x^2}} + n(1+\cosh x)^{n-1}\cosh x\)} & M1 A1 & {\msNote{M1}{Differentiates twice, the second time by the product rule, reaching terms in \((1+\cosh x)^{n-2}\sinh^2 x\) and \((1+\cosh x)^{n-1}\cosh x\). Condone coefficient slips.}\\ \msNote{A1}{Correct, unsimplified.}} \\
& {\(\sinh^2 x = (\cosh x - 1)(\cosh x + 1)\)\\[4pt] \(\Rightarrow \frac{\mathrm{d}^2y}{\mathrm{d}x^2} = n(1+\cosh x)^{n-1}\)\\[4pt] \(\phantom{\Rightarrow \frac{\mathrm{d}^2y}{\mathrm{d}x^2} =} \times\left((n-1)(\cosh x - 1) + \cosh x\right)\)} & M1 & {Replaces \(\sinh^2 x\) by \(\cosh^2 x - 1\) oe and takes out the factor \((1+\cosh x)^{n-1}\). Condone sign slips.} \\
& {\(= n(1+\cosh x)^{n-1}\left(n(1+\cosh x) - (2n-1)\right)\)\\[4pt] \(= n^2(1+\cosh x)^n - n(2n-1)(1+\cosh x)^{n-1}\)} & A1~AG & {Printed result reached without error from a correct factorised form; combined steps accepted.} \\
\nopagebreak & & [4] & \\
17(a)(ii) & {\(\frac{\mathrm{d}^2}{\mathrm{d}x^2}(1+\cosh x)^{n-1} = (n-1)^2(1+\cosh x)^{n-1}\)\\[4pt] \(\phantom{\frac{\mathrm{d}^2}{\mathrm{d}x^2}(1+\cosh x)^{n-1}} - (n-1)(2n-3)(1+\cosh x)^{n-2}\)} & M1 & {Applies (a)(i) to both terms, replacing \(n\) by \(n-1\) for the second, or differentiates twice by another valid method. Condone coefficient slips.} \\
& {\(\frac{\mathrm{d}^4y}{\mathrm{d}x^4} = n^4(1+\cosh x)^n\)\\[4pt] \(\phantom{\frac{\mathrm{d}^4y}{\mathrm{d}x^4} =} - n(2n-1)(2n^2-2n+1)(1+\cosh x)^{n-1}\)\\[4pt] \(\phantom{\frac{\mathrm{d}^4y}{\mathrm{d}x^4} =} + n(n-1)(2n-1)(2n-3)(1+\cosh x)^{n-2}\)} & A1 & {Any correct equivalent; need not be collected.} \\
\nopagebreak & & [2] & \\
17(b) & {\(y(0) = 2^n\)\\[4pt] \(y''(0) = n^2 2^n - n(2n-1)2^{n-1} = \ldots\)\\[4pt] \(y^{(4)}(0) = n^4 2^n - n(2n-1)(2n^2-2n+1)2^{n-1}\)\\[4pt] \(\phantom{y^{(4)}(0) =} + n(n-1)(2n-1)(2n-3)2^{n-2} = \ldots\)} & M1 & {Substitutes \(\cosh 0 = 1\) into \(y\) and their (a)(i) and (a)(ii) expressions in powers of \((1+\cosh x)\).} \\
& {\(y''(0) = n2^{n-1}\), \(y^{(4)}(0) = n(3n-1)2^{n-2}\)} & A1 & {Both correct, oe.} \\
& {\(y = 2^n + \frac{x^2}{2!}\text{"}n2^{n-1}\text{"} + \frac{x^4}{4!}\text{"}n(3n-1)2^{n-2}\text{"} + \cdots\)} & dM1 & {Uses the Maclaurin formula with their values and \(2!\), \(4!\), with no odd-power terms. No justification that the odd derivatives vanish is needed.} \\
& {\(y = 2^n + n2^{n-2}x^2 + \frac{n(3n-1)2^{n-5}}{3}x^4 + \cdots\)} & A1 & {All three terms, with \(x^2\) coefficient \(n2^{n-2}\) or \(\frac{n2^n}{4}\), and \(x^4\) coefficient \(\frac{n(3n-1)2^{n-5}}{3}\) or \(\frac{n(3n-1)2^n}{96}\) only. A0 for any other form, such as \(\frac{n(3n-1)2^{n-2}}{24}\).} \\
\nopagebreak & & [4] & \\
\nopagebreak \SetCell[c=2]{r} Total for Question 17 & & 10 & \\
\end{msscheme}

\end{document}
